\documentclass[mathserif, handout]{beamer}

\usepackage{url}
\input{includes}
\input{tikz_figures}

\usetikzlibrary{shapes.gates.logic.US}

\newcommand{\ttab}[1]{\texttt{#1}}

% ------------------------------------------------------------
% seven-segment figures
% ------------------------------------------------------------
% A digit is drawn in a 1 x 2 box with its lower-left corner at the origin of
% the current scope. Segments are named the standard way:
%       a
%     f   b
%       g
%     e   c
%       d
% Drawn like a real LED module: a black face, dark red segments that are
% off, and bright red segments that are on. Each segment is a pointed bar, so
% neighbouring segments meet at a mitre with a small gap between them.
\definecolor{seglit}{RGB}{255,48,36}     % on
\definecolor{seghot}{RGB}{255,196,40}    % on, and the segment being discussed
\definecolor{segoff}{RGB}{48,14,12}      % off
\definecolor{segface}{RGB}{16,16,16}     % the module's face
\newcommand{\segt}{0.18}                 % bar thickness, in digit units
% name / orientation / start / end, in a 1 x 2 box
\newcommand{\segcoords}{%
  a/h/0.09/2/0.91/2,     b/v/1/1.91/1/1.09,     c/v/1/0.91/1/0.09,
  d/h/0.09/0/0.91/0,     e/v/0/0.91/0/0.09,     f/v/0/1.91/0/1.09,
  g/h/0.09/1/0.91/1}
\newcommand{\segbar}[6]{% colour, h|v, x1, y1, x2, y2
  \if#2h
    \fill[#1] (#3,#4) -- ++(\segt/2,\segt/2) -- (#5-\segt/2,#6+\segt/2) -- (#5,#6)
              -- (#5-\segt/2,#6-\segt/2) -- (#3+\segt/2,#4-\segt/2) -- cycle;
  \else
    \fill[#1] (#3,#4) -- ++(\segt/2,-\segt/2) -- (#5+\segt/2,#6+\segt/2) -- (#5,#6)
              -- (#5-\segt/2,#6+\segt/2) -- (#3-\segt/2,#4-\segt/2) -- cycle;
  \fi}

% \sevenseg{position}{scale}{lit segments}{emphasised segments}
% Lit segments are always full red, so every digit reads correctly. An
% emphasised segment is drawn amber instead, but only where it is lit; listed
% on its own (no lit segments) it is drawn amber to point it out.
\newcommand{\sevenseg}[4]{%
  % The lists may arrive as macros (\segFive), so expand them before looping.
  \edef\seglit{#3}\edef\seghot{#4}%
  \begin{scope}[shift={#1}, scale=#2]
    \fill[segface, rounded corners=0.12cm*#2] (-0.36,-0.34) rectangle (1.36,2.34);
    \foreach \s/\o/\x/\y/\u/\v in \segcoords {
      \segbar{segoff}{\o}{\x}{\y}{\u}{\v}
      \foreach \t in \seglit {%
        \edef\sega{\s}\edef\segb{\t}%
        \ifx\sega\segb \segbar{seglit}{\o}{\x}{\y}{\u}{\v}\fi}%
      \foreach \t in \seghot {%
        \edef\sega{\s}\edef\segb{\t}%
        \ifx\sega\segb
          % amber if this digit lights it, or if no digit is given at all
          \def\segshow{0}%
          \ifx\seglit\empty \def\segshow{1}\fi
          \foreach \w in \seglit {\edef\segc{\w}\ifx\segc\segb \xdef\segshow{1}\fi}%
          \if1\segshow \segbar{seghot}{\o}{\x}{\y}{\u}{\v}\fi
        \fi}%
    }
  \end{scope}}

% The segments lit for each hex digit.
\newcommand{\segZero}{a,b,c,d,e,f}
\newcommand{\segOne}{b,c}
\newcommand{\segTwo}{a,b,d,e,g}
\newcommand{\segThree}{a,b,c,d,g}
\newcommand{\segFour}{b,c,f,g}
\newcommand{\segFive}{a,c,d,f,g}
\newcommand{\segSix}{a,c,d,e,f,g}
\newcommand{\segSeven}{a,b,c}
\newcommand{\segEight}{a,b,c,d,e,f,g}
\newcommand{\segNine}{a,b,c,d,f,g}
\newcommand{\segA}{a,b,c,e,f,g}
\newcommand{\segB}{c,d,e,f,g}
\newcommand{\segC}{a,d,e,f}
\newcommand{\segD}{b,c,d,e,g}
\newcommand{\segE}{a,d,e,f,g}
\newcommand{\segF}{a,e,f,g}

% A small inline digit for tables: \segglyph{lit}{highlighted}
\newcommand{\segglyph}[2]{%
  \tikz[baseline=-0.1ex]{\sevenseg{(0,0)}{0.19}{#1}{#2}}}

% a row of the segment table: \segrow{input bits}{glyph}{output}
\newcommand{\segrow}[3]{\ttab{#1} & #2 & \ttab{#3} \\}

% title
\title{Logic Devices}
\subtitle{Computer Organization}
\author[C. Gummaluru]{Chandra Gummaluru}
\titlegraphic{\includegraphics[width=0.15\linewidth]{images/signature.pdf}}
\institute[]{\url{chandra-gummaluru.github.io} }
\date[Computer Organization]{}

\begin{document}

\begin{frame}
  \maketitle
\end{frame}

\begin{frame}
  \begin{center}
    \highlight{\textbf{logic device}} \\[2mm]
    \graytext{a reusable block, built from gates, \\ that performs one common task}
  \end{center}
\end{frame}

\begin{frame}
  We treat each device as a single symbol and reason about what it does, not how its gates are wired.\\[3mm]
  \begin{itemize}
    \item \highlight{\textbf{multiplexer}}\\
    \graytext{selects one input to pass through} \\[2mm]
    \item \highlight{\textbf{demultiplexer}}\\
    \graytext{selects one output to pass the input through} \\[2mm]
    \item \highlight{\textbf{decoder}}\\
    \graytext{turns a number into a one-hot signal} \\[2mm]
    \item \highlight{\textbf{encoder}}\\
    \graytext{turns a one-hot signal back into a number}
  \end{itemize}
  \vspace{3mm}
These devices give us the vocabulary to describe larger circuits without drawing every gate.
\end{frame}

% ------------------------------------------------------------
% multiplexer
% ------------------------------------------------------------
\begin{frame}{Multiplexers}
  \begin{center}
    \Large \highlight{\textbf{multiplexers}}
  \end{center}
\end{frame}

\begin{frame}
  A \highlight{\textbf{multiplexer}} passes one of its data inputs to the output. A select input chooses which one.\\[3mm]
  \begin{minipage}{0.55\textwidth}
    \centering
    \begin{tikzpicture}
      \draw[thick, fill=gray!30!white!20] (0,1.3) -- (0,-1.3) -- (1.5,-0.7) -- (1.5,0.7) -- cycle;
      \draw[thick] (0,0.6) -- ++(-0.7,0) node[left, smallgraytext] {$\ttab{d}_{\ttab{0}}$};
      \draw[thick] (0,-0.6) -- ++(-0.7,0) node[left, smallgraytext] {$\ttab{d}_{\ttab{1}}$};
      \draw[thick] (1.5,0) -- ++(0.7,0) node[right, smallgraytext] {\ttab{y}};
      \draw[thick] (0.75,-1) -- ++(0,-0.6) node[below, smallgraytext] {\ttab{s}};
    \end{tikzpicture}
  \end{minipage}%
  \begin{minipage}{0.45\textwidth}
    \centering
    \renewcommand{\arraystretch}{1.3}
    \begin{tabular}{c|c}
      \ttab{s} & \ttab{y} \\ \hline
      \ttab{0} & $\ttab{d}_{\ttab{0}}$ \\
      \ttab{1} & $\ttab{d}_{\ttab{1}}$ \\
    \end{tabular}
  \end{minipage}
\end{frame}

\begin{frame}
  Inside, the select line either enables one \texttt{AND} gate or the other. Their outputs are combined by an \texttt{OR}.\\[1mm]
  \begin{figure}
    \centering
    % needs: \usetikzlibrary{shapes.gates.logic.US, calc, positioning}
    \begin{tikzpicture}[scale=0.95, transform shape]
      \node[and gate US, draw, thick, scale=1.2] (g0) at (0,1.6)  {};
      \node[and gate US, draw, thick, scale=1.2] (g1) at (0,-1.6) {};
      \node[or  gate US, draw, thick, scale=1.2] (o)  at (2.6,0)  {};
      \node[not gate US, draw, thick, scale=0.9] (n)  at (-3.0,0) {};

      % data inputs
      \draw[thick] (g0.input 1) -- ++(-0.6,0) node[left, smallgraytext] {$\ttab{d}_{\ttab{0}}$};
      \draw[thick] (g1.input 1) -- ++(-0.6,0) node[left, smallgraytext] {$\ttab{d}_{\ttab{1}}$};

      % NOT output feeds g0's second input
      \draw[thick] (n.output) |- (g0.input 2);

      % s: one horizontal run into the NOT gate, tapped down to g1's second input
      \coordinate (sj) at (-4.2,0);
      \draw[thick] (sj) -- (n.input);
      \draw[thick] (sj) -- ++(-0.6,0) node[left, smallgraytext] {\ttab{s}};
      \draw[thick] (sj) |- (g1.input 2);
      \fill (sj) circle (1.4pt);

      % AND outputs into the OR: drop/rise just before the gate, then enter from the left
      \draw[thick] (g0.output) -| ([xshift=-4mm]o.input 1) -- (o.input 1);
      \draw[thick] (g1.output) -| ([xshift=-4mm]o.input 2) -- (o.input 2);

      \draw[thick] (o.output) -- ++(0.5,0) node[right, smallgraytext] {\ttab{y}};
    \end{tikzpicture}
  \end{figure}
  \begin{center}
    \graytext{$\ttab{y} = \overline{\ttab{s}}\cdot\ttab{d}_{\ttab{0}} + \ttab{s}\cdot\ttab{d}_{\ttab{1}}$}
  \end{center}
\end{frame}

\begin{frame}
  With more select inputs, a multiplexer chooses among more data inputs.\\[3mm]
  \begin{center}
    \begin{tikzpicture}
      \draw[thick, fill=gray!30!white!20] (0,1.8) -- (0,-1.8) -- (1.6,-1.0) -- (1.6,1.0) -- cycle;
      \foreach \y/\l in {1.2/$\ttab{d}_{\ttab{0}}$, 0.4/$\ttab{d}_{\ttab{1}}$, -0.4/$\ttab{d}_{\ttab{2}}$, -1.2/$\ttab{d}_{\ttab{3}}$}
        \draw[thick] (0,\y) -- ++(-0.7,0) node[left, smallgraytext] {\ttab{\l}};
      \draw[thick] (1.6,0) -- ++(0.7,0) node[right, smallgraytext] {\ttab{y}};
      \draw[thick] (0.55,-1.525) -- (0.55,-2.0) node[below, smallgraytext] {$\ttab{s}_{\ttab{0}}$};
      \draw[thick] (1.0,-1.3) -- (1.0,-2.0) node[below, smallgraytext] {$\ttab{s}_{\ttab{1}}$};
    \end{tikzpicture}
  \end{center}
  \begin{center}
    \graytext{the select inputs read as a number that picks the data input}
  \end{center}
\end{frame}
\begin{frame}{De-Multiplexers}
  \begin{center}
    \Large \highlight{\textbf{de-multiplexers}}
  \end{center}
\end{frame}

\begin{frame}
  A \highlight{\textbf{de-multiplexer}} does the reverse. It sends one data input to the one output chosen by the select line.\\[3mm]
  \begin{center}
    \begin{tikzpicture}
      % mirror-image trapezoid (narrow on the left)
      \draw[thick, fill=gray!30!white!20] (1.5,1.3) -- (1.5,-1.3) -- (0,-0.7) -- (0,0.7) -- cycle;
      \draw[thick] (0,0) -- ++(-0.7,0) node[left, smallgraytext] {\ttab{d}};
      \draw[thick] (1.5,0.6) -- ++(0.7,0) node[right, smallgraytext] {$\ttab{y}_{\ttab{0}}$};
      \draw[thick] (1.5,-0.6) -- ++(0.7,0) node[right, smallgraytext] {$\ttab{y}_{\ttab{1}}$};
      \draw[thick] (0.75,-1) -- ++(0,-0.6) node[below, smallgraytext] {\ttab{s}};
    \end{tikzpicture}
  \end{center}
  \begin{center}
    \graytext{a multiplexer merges many lines into one, a demultiplexer splits one into many}
  \end{center}
\end{frame}

\begin{frame}{Decoder}
  \begin{center}
    \Large \highlight{\textbf{decoders}}
  \end{center}
\end{frame}

\begin{frame}
  A \highlight{\textbf{decoder}} reads its inputs as a number and sets exactly one output to \ttab{1}.\\[3mm]
  \begin{minipage}{0.45\textwidth}
    \centering
    \begin{tikzpicture}
      \draw[thick, fill=gray!30!white!20] (0,-1.6) rectangle (1.6,1.6);
      \draw[thick] (0,0.6) -- ++(-0.7,0) node[left, smallgraytext] {$\ttab{a}_{\ttab{1}}$};
      \draw[thick] (0,-0.6) -- ++(-0.7,0) node[left, smallgraytext] {$\ttab{a}_{\ttab{0}}$};
      \foreach \y/\l in {1.2/$\ttab{y}_{\ttab{3}}$, 0.4/$\ttab{y}_{\ttab{2}}$, -0.4/$\ttab{y}_{\ttab{1}}$, -1.2/$\ttab{y}_{\ttab{0}}$}
        \draw[thick] (1.6,\y) -- ++(0.7,0) node[right, smallgraytext] {\ttab{\l}};
    \end{tikzpicture}
  \end{minipage}%
  \begin{minipage}{0.55\textwidth}
    \centering
    \renewcommand{\arraystretch}{1.2}
    \begin{tabular}{cc|cccc}
      $\ttab{a}_{\ttab{1}}$ & $\ttab{a}_{\ttab{0}}$ & $\ttab{y}_{\ttab{3}}$ & $\ttab{y}_{\ttab{2}}$ & $\ttab{y}_{\ttab{1}}$ & $\ttab{y}_{\ttab{0}}$ \\ \hline
      \ttab{0} & \ttab{0} & \ttab{0} & \ttab{0} & \ttab{0} & \ttab{1} \\
      \ttab{0} & \ttab{1} & \ttab{0} & \ttab{0} & \ttab{1} & \ttab{0} \\
      \ttab{1} & \ttab{0} & \ttab{0} & \ttab{1} & \ttab{0} & \ttab{0} \\
      \ttab{1} & \ttab{1} & \ttab{1} & \ttab{0} & \ttab{0} & \ttab{0} \\
    \end{tabular}
  \end{minipage}
\\[3mm]
This is called a one-hot output.
\end{frame}


% ------------------------------------------------------------
% encoder
% ------------------------------------------------------------
\begin{frame}{Encoders}
  \begin{center}
    \Large \highlight{\textbf{encoders}}
  \end{center}
\end{frame}

\begin{frame}
  An \highlight{\textbf{encoder}} does the reverse of a decoder. Given a one-hot input, it outputs the number of the input that is \ttab{1}.\\[3mm]
  \begin{minipage}{0.45\textwidth}
    \centering
    \begin{tikzpicture}
      \draw[thick, fill=gray!30!white!20] (0,-1.6) rectangle (1.6,1.6);
      \foreach \y/\l in {1.2/$\ttab{d}_{\ttab{3}}$, 0.4/$\ttab{d}_{\ttab{2}}$, -0.4/$\ttab{d}_{\ttab{1}}$, -1.2/$\ttab{d}_{\ttab{0}}$}
        \draw[thick] (0,\y) -- ++(-0.7,0) node[left, smallgraytext] {\ttab{\l}};
      \draw[thick] (1.6,0.6) -- ++(0.7,0) node[right, smallgraytext] {$\ttab{a}_{\ttab{1}}$};
      \draw[thick] (1.6,-0.6) -- ++(0.7,0) node[right, smallgraytext] {$\ttab{a}_{\ttab{0}}$};
    \end{tikzpicture}
  \end{minipage}%
  \begin{minipage}{0.55\textwidth}
    \centering
    \renewcommand{\arraystretch}{1.2}
    \begin{tabular}{cccc|cc}
      $\ttab{d}_{\ttab{3}}$ & $\ttab{d}_{\ttab{2}}$ & $\ttab{d}_{\ttab{1}}$ & $\ttab{d}_{\ttab{0}}$ & $\ttab{a}_{\ttab{1}}$ & $\ttab{a}_{\ttab{0}}$ \\ \hline
      \ttab{0} & \ttab{0} & \ttab{0} & \ttab{1} & \ttab{0} & \ttab{0} \\
      \ttab{0} & \ttab{0} & \ttab{1} & \ttab{0} & \ttab{0} & \ttab{1} \\
      \ttab{0} & \ttab{1} & \ttab{0} & \ttab{0} & \ttab{1} & \ttab{0} \\
      \ttab{1} & \ttab{0} & \ttab{0} & \ttab{0} & \ttab{1} & \ttab{1} \\
    \end{tabular}
  \end{minipage}
\end{frame}

% ------------------------------------------------------------
% don't-care values
% ------------------------------------------------------------
\begin{frame}{Don't-Care Values}
  \begin{center}
    \Large \highlight{\textbf{don't-care values}}
  \end{center}
\end{frame}

\begin{frame}
  Real devices, like a decoder for a digit display, have inputs that never occur. We mark their output as a \highlight{\textbf{don't-care}}, written \ttab{X}.\\[4mm]
  \begin{center}
    \renewcommand{\arraystretch}{1.2}
    \begin{tabular}{cc|c}
      $\ttab{a}_{\ttab{1}}$ & $\ttab{a}_{\ttab{0}}$ & \ttab{y} \\ \hline
      \ttab{0} & \ttab{0} & \ttab{1} \\
      \ttab{0} & \ttab{1} & \ttab{0} \\
      \ttab{1} & \ttab{0} & \ttab{X} \\
      \ttab{1} & \ttab{1} & \ttab{X} \\
    \end{tabular}
  \end{center}
    When we simplify, we may read each \ttab{X} as whichever value makes the groups larger. A don't-care is a free choice, used only when it helps.
\end{frame}

% ------------------------------------------------------------
% seven-segment decoder
% ------------------------------------------------------------
\begin{frame}{Seven-Segment Decoders}
  \begin{center}
    \Large \highlight{\textbf{seven-segment decoders}}
  \end{center}
\end{frame}

\begin{frame}
  A \highlight{\textbf{seven-segment display}} draws a digit by lighting some of seven LEDs, numbered \ttab{0} to \ttab{6}. Different patterns give all sixteen hex digits.\\[4mm]
  \begin{minipage}{0.36\textwidth}
    \centering
    \begin{tikzpicture}
      % scale 1.2: the face runs from (-0.43,-0.41) to (1.63,2.81)
      \sevenseg{(0,0)}{1.2}{\segEight}{}
      \node[smallgraytext] at (0.6, 3.05)  {\ttab{0}};
      \node[smallgraytext] at (1.9, 1.8)   {\ttab{1}};
      \node[smallgraytext] at (1.9, 0.6)   {\ttab{2}};
      \node[smallgraytext] at (0.6, -0.65) {\ttab{3}};
      \node[smallgraytext] at (-0.7, 0.6)  {\ttab{4}};
      \node[smallgraytext] at (-0.7, 1.8)  {\ttab{5}};
      % 6 is the middle bar: label it outside, with a leader through the gap
      \draw[gray!70, thin] (-0.55, 1.2) -- (0.2, 1.2);
      \node[smallgraytext, anchor=east] at (-0.55, 1.2) {\ttab{6}};
    \end{tikzpicture}
  \end{minipage}%
  \begin{minipage}{0.64\textwidth}
    \centering
    \begin{tikzpicture}
      \foreach \d/\l [count=\i from 0] in
          {\segZero/0, \segOne/1, \segTwo/2, \segThree/3,
           \segFour/4, \segFive/5, \segSix/6, \segSeven/7,
           \segEight/8, \segNine/9, \segA/A, \segB/b,
           \segC/C, \segD/d, \segE/E, \segF/F} {
        \pgfmathtruncatemacro{\cx}{mod(\i,4)}
        \pgfmathtruncatemacro{\cy}{\i/4}
        \sevenseg{({\cx*1.15}, {-\cy*1.35})}{0.32}{\d}{}
        \node[smallgraytext] at ({\cx*1.15 + 0.16}, {-\cy*1.35 - 0.3}) {\ttab{\l}};
      }
    \end{tikzpicture}
  \end{minipage}\\[4mm]
  \graytext{\ttab{b} and \ttab{d} are lower case, so they aren't mistaken for \ttab{8} and \ttab{0}}
\end{frame}

\begin{frame}
  A \highlight{\textbf{seven-segment decoder}} reads a 4-bit number and outputs the seven segment signals, $\ttab{s}_{\ttab{0}}$ to $\ttab{s}_{\ttab{6}}$, that draw it.\\[4mm]
  \begin{center}
    \begin{tikzpicture}
      % the decoder
      \draw[thick, fill=gray!30!white!20] (0,-1.9) rectangle (2.2,1.9);
      \node[align=center] at (1.1,0) {\small decoder};
      % inputs, showing 0101
      \foreach \y/\i/\v in {1.2/3/0, 0.4/2/1, -0.4/1/0, -1.2/0/1} {
        \draw[thick] (0,\y) -- ++(-0.8,0) node[left, smallgraytext] {$\ttab{x}_{\ttab{\i}}$};
        \node[font=\ttfamily\small] at (-0.4,\y+0.2) {\v};
      }
      % outputs a..g, with the values for 5
      \foreach \v [count=\k from 0] in {1, 0, 1, 1, 0, 1, 1} {
        \pgfmathsetmacro{\yy}{1.5 - \k*0.5}
        \draw[thick] (2.2,\yy) -- ++(1.2,0) node[right, smallgraytext] {$\ttab{s}_{\ttab{\k}}$};
        \node[font=\ttfamily\small] at (2.8,\yy+0.18) {\v};
      }
      % the display it drives
      \sevenseg{(4.9,-1.2)}{1.2}{\segFive}{}
    \end{tikzpicture}
  \end{center}
  \begin{center}
    \graytext{input \ttab{0101} lights segments \ttab{0}, \ttab{2}, \ttab{3}, \ttab{5} and \ttab{6}, drawing a \ttab{5}}
  \end{center}
\end{frame}

\begin{frame}
  Each segment is its own function of the four inputs. Segment \ttab{0} is lit for every digit except \ttab{1}, \ttab{4}, \ttab{b} and \ttab{d}.\\[3mm]
  \begin{minipage}{0.22\textwidth}
    \centering
    \begin{tikzpicture}
      \sevenseg{(0,0)}{0.95}{}{a}
      \node[smallgraytext] at (0.475, 2.55) {segment \ttab{0}};
    \end{tikzpicture}
  \end{minipage}%
  \begin{minipage}{0.78\textwidth}
    \centering
    \small
    \renewcommand{\arraystretch}{1.3}
    \begin{tabular}{c@{\hspace{3mm}}c@{\hspace{3mm}}c}
      \graytext{\ttab{x}} & & $\ttab{s}_{\ttab{0}}$ \\ \hline
      \segrow{0000}{\segglyph{\segZero}{a}}{1}
      \segrow{0001}{\segglyph{\segOne}{}}{0}
      \segrow{0010}{\segglyph{\segTwo}{a}}{1}
      \segrow{0011}{\segglyph{\segThree}{a}}{1}
      \segrow{0100}{\segglyph{\segFour}{}}{0}
      \segrow{0101}{\segglyph{\segFive}{a}}{1}
      \segrow{0110}{\segglyph{\segSix}{a}}{1}
      \segrow{0111}{\segglyph{\segSeven}{a}}{1}
    \end{tabular}%
    \hspace{8mm}%
    \begin{tabular}{c@{\hspace{3mm}}c@{\hspace{3mm}}c}
      \graytext{\ttab{x}} & & $\ttab{s}_{\ttab{0}}$ \\ \hline
      \segrow{1000}{\segglyph{\segEight}{a}}{1}
      \segrow{1001}{\segglyph{\segNine}{a}}{1}
      \segrow{1010}{\segglyph{\segA}{a}}{1}
      \segrow{1011}{\segglyph{\segB}{}}{0}
      \segrow{1100}{\segglyph{\segC}{a}}{1}
      \segrow{1101}{\segglyph{\segD}{}}{0}
      \segrow{1110}{\segglyph{\segE}{a}}{1}
      \segrow{1111}{\segglyph{\segF}{a}}{1}
    \end{tabular}
  \end{minipage}\\[3mm]
  \graytext{the other six segments each get a table like this one}
\end{frame}

\begin{frame}
  Segment \ttab{0} is off in only four rows, so a \highlight{\textbf{decoder}} builds it directly: we \ttab{NOR} the outputs for \ttab{1}, \ttab{4}, \ttab{11} and \ttab{13}.\\[3mm]
  \begin{center}
    \begin{tikzpicture}
      % a 4-to-16 decoder
      \draw[thick, fill=gray!30!white!20] (0,-2.2) rectangle (1.6,2.2);
      \node[align=center] at (0.8,0) {\small 4-to-16 \\[-1mm] \small decoder};
      \draw[thick] (0,0) -- ++(-0.8,0) node[left, smallgraytext] {\ttab{x}};
      \draw[thick] ($(-0.4,0)+(-0.08,-0.15)$) -- ($(-0.4,0)+(0.08,0.15)$);
      \node[smallgraytext] at ($(-0.4,0)+(0,0.32)$) {\ttab{4}};

      % sixteen outputs, y0 at the top; the four that feed the gate are dark
      \foreach \k in {0,...,15} {
        \pgfmathsetmacro{\yy}{2.0 - \k*0.2667}
        \ifnum\k=1 \coordinate (o\k) at (1.6,\yy);
        \else\ifnum\k=4 \coordinate (o\k) at (1.6,\yy);
        \else\ifnum\k=11 \coordinate (o\k) at (1.6,\yy);
        \else\ifnum\k=13 \coordinate (o\k) at (1.6,\yy);
        \else \draw[gray!50] (1.6,\yy) -- ++(0.35,0);
        \fi\fi\fi\fi
      }
      \foreach \k in {0,1,4,11,13,15} {
        \pgfmathsetmacro{\yy}{2.0 - \k*0.2667}
        \node[smallgraytext, anchor=east, font=\tiny] at (1.55,\yy) {\k};
      }

      % the NOR gate
      \node[nor gate US, draw, thick, logic gate inputs=nnnn, scale=1.3] (n) at (4.4,0) {};
      % outer pair turns furthest right, inner pair nearer, so no wires cross
      \foreach \k/\i/\x/\y in {1/1/3.3/1.733, 4/2/3.0/0.933, 11/3/3.0/-0.933, 13/4/3.3/-1.467} {
        \draw[thick] (o\k) -- (\x,\y) |- (n.input \i);
        \fill[black] (o\k) circle (1.2pt);
      }
      \draw[thick] (n.output) -- ++(0.7,0) node[right, smallgraytext] {$\ttab{s}_{\ttab{0}}$};
      \sevenseg{(7.55,-0.9)}{0.9}{}{a}
    \end{tikzpicture}
  \end{center}
  \begin{center}
    \graytext{a digit display for \ttab{0}--\ttab{9} only would make rows \ttab{1010}--\ttab{1111} don't-cares}
  \end{center}
\end{frame}

\end{document}
